%----------------------------------------------------------------------------------------
%    PACKAGES AND THEMES
%----------------------------------------------------------------------------------------
\documentclass[handout, aspectratio=169,xcolor=dvipsnames]{beamer}
\makeatletter
\def\input@path{{theme/}}
\makeatother
\usetheme{CleanEasy}
\usepackage[utf8]{inputenc}
\usepackage{lmodern}
\usepackage[T1]{fontenc}
% \usepackage[brazil]{babel}
\usepackage{fix-cm}
\usepackage{amsmath}
\usepackage{mathtools}
\usepackage{listings}
\usepackage{xcolor}
\usepackage{hyperref}
\usepackage{graphicx} % Allows including images
\usepackage{booktabs} % Allows the use of \toprule, \midrule and \bottomrule in tables
\usepackage{tikz}
\usetikzlibrary{positioning, shapes, arrows, calc, decorations.pathreplacing, arrows.meta, backgrounds, patterns, overlay-beamer-styles}
\usepackage{etoolbox}
\usepackage{animate}
\usepackage{multimedia}
\usepackage{media9}

\usepackage{subfiles}


\newcommand{\bx}{\mathbf{x}}
\newcommand{\by}{\mathbf{y}}
\newcommand{\pd}{p_\text{data}}
\newcommand{\pt}{p_\theta}
\newcommand{\nbx}{\nabla_{\bx}}

\newcommand{\btx}{\mathbf{\tilde{x}}}
\newcommand{\qs}{q_{\sigma}}
\newcommand{\nbtx}{\nabla_{\btx}}

%----------------------------------------------------------------------------------------
%    LAYOUT CONFIGURATION
%----------------------------------------------------------------------------------------
\input{configs/configs}
%----------------------------------------------------------------------------------------

% References from references.bib
\usepackage[backend=bibtex,sorting=none]{biblatex}
\addbibresource{references.bib}


\begin{document}

\section{Week 3}


\begin{frame}{Old architecture based on LorentzNet}

    $x$ = Cartesian particle 4-momenta, $\phi_\cdot$ = learnable functions, $\psi(\cdot) = \text{sgn}(\cdot) \log(|\cdot| + 1)$, $c = $ hyperparameter, $N = \#$ particles, $|| \cdot ||, \langle \cdot \rangle$ = Minkowski norm and inner product
    \vspace{0.5cm}
    \begin{columns}
        \begin{column}{0.5\textwidth}
            \textbf{LGEB (\cite{gongEfficientLorentzEquivariant2022})}
            \begin{align*}
                m_{ij}^l &= \phi_e(h_i^l, h_j^l, \psi(||x_i^l - x_j^l||), \psi(\langle x_i^l, x_j^l \rangle)) \\
                m_{ij}^l &= \phi_m(m_{ij}^l) \cdot m_{ij}^l
            \end{align*}
            \begin{align*}
                x_i^{l+1} &= x_i^l + \frac{c}{N} \sum\limits_{j \in [N]} \phi_x(m_{ij}^l) \cdot |x_i - x_j|         
            \end{align*}
            \begin{align*}
                h_i^{l+1} &= h_i^l + \phi_h(h_i^l, \sum\limits_{j \in [N]} \phi_w(m_{ij}^l) m_{ij}^l)
            \end{align*}
        \end{column}
        \begin{column}{0.5\textwidth}
            \textbf{Simple LE block this week}
            \begin{align*}
                m^l_{ij} &= \phi_e(\psi(||x_i^l - x_j^l||), \psi(\langle x_i^l, x_j^l \rangle))\\
                m^l_{ij} &= \phi_m(m^l_{ij}) \cdot m^l_{ij}
            \end{align*}
            \begin{align*}
                x_i^{l+1} &= x_i^l + \frac{c}{N} \sum\limits_{j \in [N]} \phi_x(m_{ij}, g) \cdot |x_j^l - x_i^l|
            \end{align*}
            $g = $ global feature vector (currently time embedding $\phi_t(t)$)
        \end{column}
    \end{columns}
\end{frame}

\begin{frame}[allowframebreaks]{New architecture}
    \begin{enumerate}
        \item Embed scalar time value to random Fourier features through two fully connected layers with 32 and 64 nodes (following \cite{wuFastPointCloud2023})
        \item Concatenate time embedding with jet conditioning information, which consists of: \begin{enumerate}
            \item One-hot encoded jet type
            \item Jet mass
            \item Jet transverse momentum
            \item Jet number of constituents
        \end{enumerate}
        \item Pass this through a fully connected layer of size 64 to produce the initial global embeddings $g_0$
        \item Concatenate $g_0$ with features $x_1^0 \dots x_n^0$ initialized from the initial distribution $p_0$.
        \item Now, we use 6 Lorentz-group equivariant layers. \begin{enumerate}
            \item The input to the $l+1$-th layer is $g^l$ and $x^{l}_n \dots x^l_N$.
            \item Calculate messages $m^l_{ij} = \phi_m(\psi(||x_i^l - x_j^l||), \psi(\langle x_i^l, x_j^l \rangle))$
            \item Update the global embedding $g_{l+1} = \phi_{eg}(g_l, \frac{c}{N} \sum\limits_{i, j \in [N]} \phi_e(m_{ij}) \bigoplus \sum\limits_{i, j \in [N]} \phi_e(m_{ij}))$
            \item Feed in the updated global embedding as well as the original embedding to update the graph embeddings $$
            x_i^{l+1} = x_i^{l} + \phi_x(g_0 \bigoplus g^{l+1} \bigoplus \frac{c}{N} \sum\limits_{i, j \in [N]} \phi_{ex}(m_{ij}) \bigoplus \sum\limits_{i, j \in [N]} \phi_{ex}(m_{ij}))
            $$
            \item Pass the results $g^{l+1}$ and $x_i^{l+1}$ to the next layer
        \end{enumerate}
        \item After the final LE layer, run the particle representations through a final round of message-passing with the final global embedding $$
        x_i^{L+1} = x_i^{L} + \phi_x(g_0 \bigoplus g^{L} \bigoplus \frac{c}{N} \sum\limits_{i, j \in [N]} \phi_{ex}(m_{ij}) \bigoplus \sum\limits_{i, j \in [N]} \phi_{ex}(m_{ij}))
        $$
    \end{enumerate}
\end{frame}



\begin{frame}[allowframebreaks]{References}
\printbibliography
\end{frame}
\end{document}
